\documentclass[11pt,a4paper]{article}
\usepackage[utf8]{inputenc}
\usepackage[T1]{fontenc}
\usepackage[brazil]{babel}
\usepackage{geometry}
\usepackage{graphicx}
\usepackage{booktabs}
\usepackage{longtable}
\usepackage{array}
\usepackage{float}
\usepackage{xcolor}
\usepackage{listings}
\usepackage{pdflscape}
\usepackage{placeins}
\usepackage[hidelinks]{hyperref}
\geometry{margin=2.1cm}
\setlength{\parskip}{0.5em}
\setlength{\parindent}{0pt}
\lstset{
  basicstyle=\ttfamily\footnotesize,
  breaklines=true,
  columns=fullflexible,
  frame=single,
  showstringspaces=false,
  keywordstyle=\color{blue},
  commentstyle=\color{gray}
}
\title{Catálogo explicado das queries de benchmark Druid}
\author{CCO DataOps}
\date{30 de junho de 2026}
\begin{document}
\maketitle


\section{Como o catálogo está organizado}

Cada query é documentada por família de carga. A família explica o conceito experimental; a query explica como esse conceito foi materializado em SQL. O objetivo não é só listar SQL, mas justificar por que a consulta mede uma propriedade relevante do Druid.

\section{Resumo das queries}

\begin{landscape}
{
\scriptsize
\setlength{\tabcolsep}{4pt}
\renewcommand{\arraystretch}{1.15}
\begin{longtable}{p{0.27\textwidth}p{0.16\textwidth}p{0.14\textwidth}p{0.21\textwidth}p{0.20\textwidth}}
\caption{Catalogo resumido das queries executadas.}\\
\toprule
Query & Familia & Padrao & Datasource & Objetivo \\
\midrule
\endfirsthead
\toprule
Query & Familia & Padrao & Datasource & Objetivo \\
\midrule
\endhead
01\_latency\_simple/00\_count\_hospitais & Latencia simples & count\_scan & dim\_hospitais\_cnes & medir latencia minima em uma dimensao grande \\
01\_latency\_simple/01\_count\_dimensoes\_catalogo & Latencia simples & union\_counts & todas\_dimensoes & medir latencia de inventario com UNION ALL entre dimensoes \\
01\_latency\_simple/02\_filtro\_hospitais\_uf & Latencia simples & filtered\_groupby & dim\_hospitais\_cnes & medir filtro seletivo em coluna de baixa cardinalidade \\
01\_latency\_simple/03\_lookup\_municipios\_mg & Latencia simples & lookup\_filter & dim\_municipios\_coordenadas & medir consulta de lookup territorial usada como filtro de painel \\
02\_olap\_pesado/00\_hospitais\_groupby\_multidimensional & OLAP pesado & groupby\_distinct & dim\_hospitais\_cnes & forcar agregacao multidimensional e count distinct na maior dimensao \\
02\_olap\_pesado/01\_terras\_indigenas\_area\_regularizacao & OLAP pesado & numeric\_aggregations & dim\_terras\_indigenas & medir agregacoes numericas em dimensao espacial com path\_json pesado \\
02\_olap\_pesado/02\_rodovias\_status\_pavimento & OLAP pesado & categorical\_groupby & dim\_rodovias\_federais & medir agrupamento em dimensao linear com varias classificacoes \\
02\_olap\_pesado/03\_ons\_linhas\_tensao\_extensao & OLAP pesado & numeric\_groupby & dim\_ons\_linhas\_transmissao & medir agregacoes ONS por tensao, extensao e criticidade \\
03\_joins\_dimensionais/00\_hospitais\_municipios\_join & Joins dimensionais & left\_join\_groupby & dim\_hospitais\_cnes, dim\_municipios\_coordenadas & simular painel de defesa civil cruzando hospitais e municipios \\
03\_joins\_dimensionais/01\_hospitais\_terras\_indigenas\_join & Joins dimensionais & nearest\_fk\_join & dim\_hospitais\_cnes, dim\_terras\_indigenas & simular painel de saude cruzando hospitais com TI mais proxima \\
03\_joins\_dimensionais/02\_ti\_municipio\_principal\_join & Joins dimensionais & dimension\_lookup\_join & dim\_terras\_indigenas, dim\_municipios\_coordenadas & simular filtro territorial de TIs por municipio principal \\
03\_joins\_dimensionais/03\_municipio\_saude\_ti\_join & Joins dimensionais & multi\_join\_aggregate & dim\_municipios\_coordenadas, dim\_hospitais\_cnes, dim\_terras\_indigenas & simular dashboard combinando municipios, hospitais e TIs \\
04\_payload\_espacial/00\_path\_json\_municipios & Payload espacial & string\_payload\_scan & dim\_municipios\_coordenadas & medir custo de ler geometrias serializadas de municipios \\
04\_payload\_espacial/01\_path\_json\_terras\_indigenas & Payload espacial & string\_payload\_scan & dim\_terras\_indigenas & medir custo de ler geometrias serializadas de TIs \\
04\_payload\_espacial/02\_path\_json\_rodovias & Payload espacial & string\_payload\_scan & dim\_rodovias\_federais & medir custo de ler geometrias serializadas de rodovias \\
04\_payload\_espacial/03\_preview\_hospitais\_path\_json & Payload espacial & result\_payload & dim\_hospitais\_cnes & simular mapa de pontos retornando payload de geometria ao cliente \\
05\_dashboard\_misto/00\_dashboard\_defesa\_civil & Dashboard misto & dashboard\_composite & dim\_municipios\_coordenadas, dim\_hospitais\_cnes, dim\_terras\_indigenas & simular painel territorial combinando municipios, hospitais e TIs \\
05\_dashboard\_misto/01\_dashboard\_energia\_logistica & Dashboard misto & union\_aggregate & dim\_ons\_linhas\_transmissao, dim\_ons\_subestacoes, dim\_rodovias\_federais & simular comparativo operacional entre energia e logistica \\
05\_dashboard\_misto/02\_dashboard\_qualidade\_operacional & Dashboard misto & quality\_dashboard & dim\_hospitais\_cnes, dim\_terras\_indigenas, dim\_municipios\_coordenadas & simular painel de governanca com qualidade geografica e relacionamentos \\
06\_fatos\_streaming/00\_count\_eventos\_total & Fatos e Streaming & count\_scan & fato\_eventos\_incendio & medir latencia baseline em tabela fato (158K linhas) \\
06\_fatos\_streaming/01\_eventos\_groupby\_bioma\_mes & Fatos e Streaming & groupby\_temporal & fato\_eventos\_incendio & agregacao bioma x mes - teste do motor GroupBy em fato com 158K linhas \\
06\_fatos\_streaming/02\_eventos\_top\_frp\_ordenado & Fatos e Streaming & groupby\_ordered\_limit & fato\_eventos\_incendio & simular o padrao exato dos Virtual Datasets (GROUP BY + ORDER BY + LIMIT) \\
06\_fatos\_streaming/03\_focos\_scan\_temporal & Fatos e Streaming & scan\_temporal & fato\_focos\_incendio & medir motor Scan em tabela de 3.3M linhas com filtro temporal \\
06\_fatos\_streaming/04\_focos\_lookup\_municipio & Fatos e Streaming & lookup\_heavy & fato\_focos\_incendio & medir custo de LOOKUP em tabela de 3.3M linhas \\
07\_virtual\_datasets/00\_ds\_eventos\_geral & Virtual Datasets & virtual\_dataset & fato\_eventos\_incendio & simular query real do dashboard de eventos geral (Virtual Dataset) \\
07\_virtual\_datasets/01\_ds\_focos\_individuais & Virtual Datasets & virtual\_dataset & fato\_focos\_incendio & simular query real do dashboard de focos individuais (Virtual Dataset) \\
07\_virtual\_datasets/02\_ds\_eventos\_energia & Virtual Datasets & virtual\_dataset & fato\_eventos\_incendio & simular query real do dashboard de eventos energia (Virtual Dataset) \\
07\_virtual\_datasets/03\_ds\_eventos\_indigenas & Virtual Datasets & virtual\_dataset & fato\_eventos\_incendio & simular query real do dashboard de eventos indigenas (Virtual Dataset) \\
07\_virtual\_datasets/04\_ds\_eventos\_logistica & Virtual Datasets & virtual\_dataset & fato\_eventos\_incendio & simular query real do dashboard de eventos logistica (Virtual Dataset) \\
\bottomrule
\end{longtable}
}
\end{landscape}


\clearpage

\section{Latencia simples}

\textbf{Conceito base.} Consultas pequenas usadas para medir o piso de resposta do Broker e a leitura basica de segmentos.

\textbf{Hipótese testada.} Se ate essas consultas oscilarem, o problema esta no cluster ou na disponibilidade dos segmentos, nao na complexidade analitica.

\textbf{Por que existe.} Funciona como controle experimental para comparar as consultas mais pesadas.

\textbf{Impacto esperado no Druid.} Pressiona pouco o Historical; serve para medir overhead de SQL, roteamento e filtros simples.

\subsection{01\_latency\_simple/00\_count\_hospitais}

\textbf{Objetivo local.} medir latencia minima em uma dimensao grande\\

\textbf{Datasource(s).} dim\_hospitais\_cnes\\

\textbf{Padrão de carga.} count\_scan\\

\textbf{Visualização Superset simulada.} indicador\\

\textbf{Motivo científico.} Funciona como controle experimental para comparar as consultas mais pesadas. Objetivo local: medir latencia minima em uma dimensao grande\\

\textbf{Impacto esperado.} Pressiona pouco o Historical; serve para medir overhead de SQL, roteamento e filtros simples.

\begin{lstlisting}[language=SQL]

SELECT COUNT(*) AS total_hospitais
FROM dim_hospitais_cnes;

\end{lstlisting}

\subsection{01\_latency\_simple/01\_count\_dimensoes\_catalogo}

\textbf{Objetivo local.} medir latencia de inventario com UNION ALL entre dimensoes\\

\textbf{Datasource(s).} todas\_dimensoes\\

\textbf{Padrão de carga.} union\_counts\\

\textbf{Visualização Superset simulada.} tabela\\

\textbf{Motivo científico.} Funciona como controle experimental para comparar as consultas mais pesadas. Objetivo local: medir latencia de inventario com UNION ALL entre dimensoes\\

\textbf{Impacto esperado.} Pressiona pouco o Historical; serve para medir overhead de SQL, roteamento e filtros simples.

\begin{lstlisting}[language=SQL]

SELECT 'municipios' AS datasource, COUNT(*) AS total FROM dim_municipios_coordenadas
UNION ALL
SELECT 'hospitais' AS datasource, COUNT(*) AS total FROM dim_hospitais_cnes
UNION ALL
SELECT 'terras_indigenas' AS datasource, COUNT(*) AS total FROM dim_terras_indigenas
UNION ALL
SELECT 'rodovias' AS datasource, COUNT(*) AS total FROM dim_rodovias_federais
UNION ALL
SELECT 'ons_linhas' AS datasource, COUNT(*) AS total FROM dim_ons_linhas_transmissao
UNION ALL
SELECT 'ons_subestacoes' AS datasource, COUNT(*) AS total FROM dim_ons_subestacoes
UNION ALL
SELECT 'ons_sistemas_isolados' AS datasource, COUNT(*) AS total FROM dim_ons_sistemas_isolados;

\end{lstlisting}

\subsection{01\_latency\_simple/02\_filtro\_hospitais\_uf}

\textbf{Objetivo local.} medir filtro seletivo em coluna de baixa cardinalidade\\

\textbf{Datasource(s).} dim\_hospitais\_cnes\\

\textbf{Padrão de carga.} filtered\_groupby\\

\textbf{Visualização Superset simulada.} barras\\

\textbf{Motivo científico.} Funciona como controle experimental para comparar as consultas mais pesadas. Objetivo local: medir filtro seletivo em coluna de baixa cardinalidade\\

\textbf{Impacto esperado.} Pressiona pouco o Historical; serve para medir overhead de SQL, roteamento e filtros simples.

\begin{lstlisting}[language=SQL]

SELECT
  LOOKUP(codigo_ibge, 'municipios_uf') AS estado_sigla,
  COUNT(*) AS hospitais_total,
  SUM(CASE WHEN funciona_24h = 'SIM' THEN 1 ELSE 0 END) AS hospitais_24h
FROM dim_hospitais_cnes
WHERE LOOKUP(codigo_ibge, 'municipios_uf') IN ('MG', 'SP', 'RJ', 'PA', 'AM')
GROUP BY 1;

\end{lstlisting}

\subsection{01\_latency\_simple/03\_lookup\_municipios\_mg}

\textbf{Objetivo local.} medir consulta de lookup territorial usada como filtro de painel\\

\textbf{Datasource(s).} dim\_municipios\_coordenadas\\

\textbf{Padrão de carga.} lookup\_filter\\

\textbf{Visualização Superset simulada.} tabela\\

\textbf{Motivo científico.} Funciona como controle experimental para comparar as consultas mais pesadas. Objetivo local: medir consulta de lookup territorial usada como filtro de painel\\

\textbf{Impacto esperado.} Pressiona pouco o Historical; serve para medir overhead de SQL, roteamento e filtros simples.

\begin{lstlisting}[language=SQL]

SELECT
  codigo_ibge,
  municipios_nome,
  estado_sigla,
  regiao_brasil
FROM dim_municipios_coordenadas
WHERE estado_sigla = 'MG';

\end{lstlisting}

\clearpage

\section{OLAP pesado}

\textbf{Conceito base.} Agregacoes, agrupamentos multidimensionais, soma, media e count distinct sobre dimensoes batch.

\textbf{Hipótese testada.} Quanto mais cardinalidade e agregacao numerica, maior a pressao sobre processamento e merge de resultados.

\textbf{Por que existe.} Representa o uso esperado de um banco OLAP em paineis de resumo e ranking.

\textbf{Impacto esperado no Druid.} Exige varredura de colunas, agregacao no Historical e merge no Broker.

\subsection{02\_olap\_pesado/00\_hospitais\_groupby\_multidimensional}

\textbf{Objetivo local.} forcar agregacao multidimensional e count distinct na maior dimensao\\

\textbf{Datasource(s).} dim\_hospitais\_cnes\\

\textbf{Padrão de carga.} groupby\_distinct\\

\textbf{Visualização Superset simulada.} tabela\_barras\\

\textbf{Motivo científico.} Representa o uso esperado de um banco OLAP em paineis de resumo e ranking. Objetivo local: forcar agregacao multidimensional e count distinct na maior dimensao\\

\textbf{Impacto esperado.} Exige varredura de colunas, agregacao no Historical e merge no Broker.

\begin{lstlisting}[language=SQL]

SELECT
  tipo_unidade_critica,
  categoria_unidade,
  classe_capacidade_saude,
  funciona_24h,
  COUNT(*) AS registros_total,
  COUNT(DISTINCT codigo_cnes) AS hospitais_distintos,
  SUM(CASE WHEN dentro_terra_indigena = 1 THEN 1 ELSE 0 END) AS dentro_ti
FROM dim_hospitais_cnes
GROUP BY
  tipo_unidade_critica,
  categoria_unidade,
  classe_capacidade_saude,
  funciona_24h
ORDER BY registros_total DESC;

\end{lstlisting}

\subsection{02\_olap\_pesado/01\_terras\_indigenas\_area\_regularizacao}

\textbf{Objetivo local.} medir agregacoes numericas em dimensao espacial com path\_json pesado\\

\textbf{Datasource(s).} dim\_terras\_indigenas\\

\textbf{Padrão de carga.} numeric\_aggregations\\

\textbf{Visualização Superset simulada.} barras\\

\textbf{Motivo científico.} Representa o uso esperado de um banco OLAP em paineis de resumo e ranking. Objetivo local: medir agregacoes numericas em dimensao espacial com path\_json pesado\\

\textbf{Impacto esperado.} Exige varredura de colunas, agregacao no Historical e merge no Broker.

\begin{lstlisting}[language=SQL]

SELECT
  estado_sigla,
  estagio_regularizacao,
  prioridade_funai,
  COUNT(*) AS terras_indigenas_total,
  SUM(area_hectares) AS area_hectares_total,
  AVG(qtd_municipios_interceptados) AS media_municipios_interceptados,
  AVG(percentual_area_interceptada) AS media_percentual_interceptado
FROM dim_terras_indigenas
GROUP BY
  estado_sigla,
  estagio_regularizacao,
  prioridade_funai
ORDER BY area_hectares_total DESC;

\end{lstlisting}

\subsection{02\_olap\_pesado/02\_rodovias\_status\_pavimento}

\textbf{Objetivo local.} medir agrupamento em dimensao linear com varias classificacoes\\

\textbf{Datasource(s).} dim\_rodovias\_federais\\

\textbf{Padrão de carga.} categorical\_groupby\\

\textbf{Visualização Superset simulada.} barras\_empilhadas\\

\textbf{Motivo científico.} Representa o uso esperado de um banco OLAP em paineis de resumo e ranking. Objetivo local: medir agrupamento em dimensao linear com varias classificacoes\\

\textbf{Impacto esperado.} Exige varredura de colunas, agregacao no Historical e merge no Broker.

\begin{lstlisting}[language=SQL]

SELECT
  codigo_rodovia,
  jurisdicao,
  classe_status_operacional,
  classe_situacao_fisica,
  classe_revestimento,
  COUNT(*) AS trechos_total
FROM dim_rodovias_federais
GROUP BY
  codigo_rodovia,
  jurisdicao,
  classe_status_operacional,
  classe_situacao_fisica,
  classe_revestimento
ORDER BY trechos_total DESC;

\end{lstlisting}

\subsection{02\_olap\_pesado/03\_ons\_linhas\_tensao\_extensao}

\textbf{Objetivo local.} medir agregacoes ONS por tensao, extensao e criticidade\\

\textbf{Datasource(s).} dim\_ons\_linhas\_transmissao\\

\textbf{Padrão de carga.} numeric\_groupby\\

\textbf{Visualização Superset simulada.} barras\\

\textbf{Motivo científico.} Representa o uso esperado de um banco OLAP em paineis de resumo e ranking. Objetivo local: medir agregacoes ONS por tensao, extensao e criticidade\\

\textbf{Impacto esperado.} Exige varredura de colunas, agregacao no Historical e merge no Broker.

\begin{lstlisting}[language=SQL]

SELECT
  estado_sigla_estimado,
  classe_tensao,
  classe_criticidade_linha,
  COUNT(*) AS linhas_total,
  SUM(extensao_km) AS extensao_km_total,
  AVG(tensao_kv) AS tensao_kv_media,
  MAX(tensao_kv) AS tensao_kv_maxima
FROM dim_ons_linhas_transmissao
GROUP BY
  estado_sigla_estimado,
  classe_tensao,
  classe_criticidade_linha
ORDER BY extensao_km_total DESC;

\end{lstlisting}

\clearpage

\section{Joins dimensionais}

\textbf{Conceito base.} Cruzamentos entre dimensoes ja enriquecidas pelo ETL para simular visoes analiticas territoriais.

\textbf{Hipótese testada.} Joins com hospitais, municipios e terras indigenas devem expor gargalos de planejamento, shuffle e merge.

\textbf{Por que existe.} Testa o limite pratico de usar Druid SQL para cruzamentos que o Superset pode exigir.

\textbf{Impacto esperado no Druid.} Pressiona Broker e Historical ao mesmo tempo, principalmente quando a maior dimensao participa.

\subsection{03\_joins\_dimensionais/00\_hospitais\_municipios\_join}

\textbf{Objetivo local.} simular painel de defesa civil cruzando hospitais e municipios\\

\textbf{Datasource(s).} dim\_hospitais\_cnes, dim\_municipios\_coordenadas\\

\textbf{Padrão de carga.} left\_join\_groupby\\

\textbf{Visualização Superset simulada.} mapa\_coropletico\\

\textbf{Motivo científico.} Testa o limite pratico de usar Druid SQL para cruzamentos que o Superset pode exigir. Objetivo local: simular painel de defesa civil cruzando hospitais e municipios\\

\textbf{Impacto esperado.} Pressiona Broker e Historical ao mesmo tempo, principalmente quando a maior dimensao participa.

\begin{lstlisting}[language=SQL]

SELECT
  m.regiao_brasil,
  m.estado_sigla,
  m.codigo_ibge,
  m.municipios_nome,
  COUNT(DISTINCT h.codigo_cnes) AS hospitais_total,
  SUM(CASE WHEN h.funciona_24h = 'SIM' THEN 1 ELSE 0 END) AS hospitais_24h
FROM dim_municipios_coordenadas m
LEFT JOIN dim_hospitais_cnes h
  ON h.codigo_ibge = m.codigo_ibge
GROUP BY
  m.regiao_brasil,
  m.estado_sigla,
  m.codigo_ibge,
  m.municipios_nome
ORDER BY hospitais_total DESC;

\end{lstlisting}

\subsection{03\_joins\_dimensionais/01\_hospitais\_terras\_indigenas\_join}

\textbf{Objetivo local.} simular painel de saude cruzando hospitais com TI mais proxima\\

\textbf{Datasource(s).} dim\_hospitais\_cnes, dim\_terras\_indigenas\\

\textbf{Padrão de carga.} nearest\_fk\_join\\

\textbf{Visualização Superset simulada.} tabela\_mapa\\

\textbf{Motivo científico.} Testa o limite pratico de usar Druid SQL para cruzamentos que o Superset pode exigir. Objetivo local: simular painel de saude cruzando hospitais com TI mais proxima\\

\textbf{Impacto esperado.} Pressiona Broker e Historical ao mesmo tempo, principalmente quando a maior dimensao participa.

\begin{lstlisting}[language=SQL]

SELECT
  t.estado_sigla,
  t.nome_terra_indigena,
  t.prioridade_funai,
  t.estagio_regularizacao,
  COUNT(DISTINCT h.codigo_cnes) AS hospitais_relacionados,
  SUM(CASE WHEN h.dentro_terra_indigena = 1 THEN 1 ELSE 0 END) AS hospitais_dentro_ti,
  SUM(CASE WHEN h.dentro_terra_indigena = 1 THEN 1 ELSE 0 END) AS hospitais_ate_15km,
  AVG(h.distancia_ti_km) AS distancia_media_km
FROM dim_terras_indigenas t
LEFT JOIN dim_hospitais_cnes h
  ON h.codigo_ti_proxima = t.codigo_ti
GROUP BY
  t.estado_sigla,
  t.nome_terra_indigena,
  t.prioridade_funai,
  t.estagio_regularizacao
ORDER BY hospitais_relacionados DESC;

\end{lstlisting}

\subsection{03\_joins\_dimensionais/02\_ti\_municipio\_principal\_join}

\textbf{Objetivo local.} simular filtro territorial de TIs por municipio principal\\

\textbf{Datasource(s).} dim\_terras\_indigenas, dim\_municipios\_coordenadas\\

\textbf{Padrão de carga.} dimension\_lookup\_join\\

\textbf{Visualização Superset simulada.} tabela\\

\textbf{Motivo científico.} Testa o limite pratico de usar Druid SQL para cruzamentos que o Superset pode exigir. Objetivo local: simular filtro territorial de TIs por municipio principal\\

\textbf{Impacto esperado.} Pressiona Broker e Historical ao mesmo tempo, principalmente quando a maior dimensao participa.

\begin{lstlisting}[language=SQL]

SELECT
  m.regiao_brasil,
  m.estado_sigla,
  m.municipios_nome,
  t.prioridade_funai,
  t.estagio_regularizacao,
  COUNT(*) AS terras_indigenas_total,
  SUM(t.area_hectares) AS area_hectares_total,
  AVG(t.qtd_municipios_interceptados) AS media_municipios_interceptados
FROM dim_terras_indigenas t
LEFT JOIN dim_municipios_coordenadas m
  ON t.codigo_ibge_principal = m.codigo_ibge
GROUP BY
  m.regiao_brasil,
  m.estado_sigla,
  m.municipios_nome,
  t.prioridade_funai,
  t.estagio_regularizacao
ORDER BY area_hectares_total DESC;

\end{lstlisting}

\subsection{03\_joins\_dimensionais/03\_municipio\_saude\_ti\_join}

\textbf{Objetivo local.} simular dashboard combinando municipios, hospitais e TIs\\

\textbf{Datasource(s).} dim\_municipios\_coordenadas, dim\_hospitais\_cnes, dim\_terras\_indigenas\\

\textbf{Padrão de carga.} multi\_join\_aggregate\\

\textbf{Visualização Superset simulada.} tabela\_mapa\\

\textbf{Motivo científico.} Testa o limite pratico de usar Druid SQL para cruzamentos que o Superset pode exigir. Objetivo local: simular dashboard combinando municipios, hospitais e TIs\\

\textbf{Impacto esperado.} Pressiona Broker e Historical ao mesmo tempo, principalmente quando a maior dimensao participa.

\begin{lstlisting}[language=SQL]

WITH hospitais_por_municipio AS (
  SELECT
    codigo_ibge,
    COUNT(DISTINCT codigo_cnes) AS hospitais_total,
    SUM(CASE WHEN dentro_terra_indigena = 1 THEN 1 ELSE 0 END) AS hospitais_ate_15km_ti
  FROM dim_hospitais_cnes
  GROUP BY codigo_ibge
),
tis_por_municipio AS (
  SELECT
    codigo_ibge_principal,
    COUNT(*) AS terras_indigenas_total,
    SUM(area_hectares) AS area_ti_hectares
  FROM dim_terras_indigenas
  GROUP BY codigo_ibge_principal
)
SELECT
  m.regiao_brasil,
  m.estado_sigla,
  m.codigo_ibge,
  m.municipios_nome,
  COALESCE(h.hospitais_total, 0) AS hospitais_total,
  COALESCE(h.hospitais_ate_15km_ti, 0) AS hospitais_ate_15km_ti,
  COALESCE(t.terras_indigenas_total, 0) AS terras_indigenas_total,
  COALESCE(t.area_ti_hectares, 0) AS area_ti_hectares
FROM dim_municipios_coordenadas m
LEFT JOIN hospitais_por_municipio h
  ON h.codigo_ibge = m.codigo_ibge
LEFT JOIN tis_por_municipio t
  ON t.codigo_ibge_principal = m.codigo_ibge;

\end{lstlisting}

\clearpage

\section{Payload espacial}

\textbf{Conceito base.} Leituras de `path\_json` e campos de geometria serializada usados por mapas no Superset.

\textbf{Hipótese testada.} Mesmo sem calculo espacial nativo, retornar strings grandes aumenta bytes e latencia de resposta.

\textbf{Por que existe.} Mede o custo operacional do contrato visual adotado para Druid/Superset.

\textbf{Impacto esperado no Druid.} Pressiona I/O, serializacao e tamanho de resposta ao cliente.

\subsection{04\_payload\_espacial/00\_path\_json\_municipios}

\textbf{Objetivo local.} medir custo de ler geometrias serializadas de municipios\\

\textbf{Datasource(s).} dim\_municipios\_coordenadas\\

\textbf{Padrão de carga.} string\_payload\_scan\\

\textbf{Visualização Superset simulada.} indicador\\

\textbf{Motivo científico.} Mede o custo operacional do contrato visual adotado para Druid/Superset. Objetivo local: medir custo de ler geometrias serializadas de municipios\\

\textbf{Impacto esperado.} Pressiona I/O, serializacao e tamanho de resposta ao cliente.

\begin{lstlisting}[language=SQL]

SELECT
  COUNT(*) AS municipios_total,
  SUM(LENGTH(path_json)) AS bytes_path_json_estimados,
  AVG(LENGTH(path_json)) AS media_path_json,
  MAX(LENGTH(path_json)) AS maior_path_json
FROM dim_municipios_coordenadas;

\end{lstlisting}

\subsection{04\_payload\_espacial/01\_path\_json\_terras\_indigenas}

\textbf{Objetivo local.} medir custo de ler geometrias serializadas de TIs\\

\textbf{Datasource(s).} dim\_terras\_indigenas\\

\textbf{Padrão de carga.} string\_payload\_scan\\

\textbf{Visualização Superset simulada.} indicador\\

\textbf{Motivo científico.} Mede o custo operacional do contrato visual adotado para Druid/Superset. Objetivo local: medir custo de ler geometrias serializadas de TIs\\

\textbf{Impacto esperado.} Pressiona I/O, serializacao e tamanho de resposta ao cliente.

\begin{lstlisting}[language=SQL]

SELECT
  COUNT(*) AS terras_indigenas_total,
  SUM(LENGTH(path_json)) AS bytes_path_json_estimados,
  AVG(LENGTH(path_json)) AS media_path_json,
  MAX(LENGTH(path_json)) AS maior_path_json
FROM dim_terras_indigenas;

\end{lstlisting}

\subsection{04\_payload\_espacial/02\_path\_json\_rodovias}

\textbf{Objetivo local.} medir custo de ler geometrias serializadas de rodovias\\

\textbf{Datasource(s).} dim\_rodovias\_federais\\

\textbf{Padrão de carga.} string\_payload\_scan\\

\textbf{Visualização Superset simulada.} indicador\\

\textbf{Motivo científico.} Mede o custo operacional do contrato visual adotado para Druid/Superset. Objetivo local: medir custo de ler geometrias serializadas de rodovias\\

\textbf{Impacto esperado.} Pressiona I/O, serializacao e tamanho de resposta ao cliente.

\begin{lstlisting}[language=SQL]

SELECT
  COUNT(*) AS trechos_total,
  SUM(LENGTH(path_json)) AS bytes_path_json_estimados,
  AVG(LENGTH(path_json)) AS media_path_json,
  MAX(LENGTH(path_json)) AS maior_path_json
FROM dim_rodovias_federais;

\end{lstlisting}

\subsection{04\_payload\_espacial/03\_preview\_hospitais\_path\_json}

\textbf{Objetivo local.} simular mapa de pontos retornando payload de geometria ao cliente\\

\textbf{Datasource(s).} dim\_hospitais\_cnes\\

\textbf{Padrão de carga.} result\_payload\\

\textbf{Visualização Superset simulada.} deck\_gl\_scatterplot\\

\textbf{Motivo científico.} Mede o custo operacional do contrato visual adotado para Druid/Superset. Objetivo local: simular mapa de pontos retornando payload de geometria ao cliente\\

\textbf{Impacto esperado.} Pressiona I/O, serializacao e tamanho de resposta ao cliente.

\begin{lstlisting}[language=SQL]

SELECT
  codigo_cnes,
  nome_hospital,
  codigo_ibge,
  latitude,
  longitude,
  path_json
FROM dim_hospitais_cnes
WHERE latitude BETWEEN -35 AND 6
  AND longitude BETWEEN -75 AND -30
LIMIT 1000;

\end{lstlisting}

\clearpage

\section{Dashboard misto}

\textbf{Conceito base.} Consultas compostas que combinam padroes de painel real: agregacao, joins, union e regras de qualidade.

\textbf{Hipótese testada.} A mistura de padroes aproxima melhor a experiencia de uso do que queries isoladas.

\textbf{Por que existe.} Valida se o perfil suporta uma sessao analitica mais proxima do Superset.

\textbf{Impacto esperado no Druid.} Exercita planejamento SQL, subconsultas, union, filtros e agregacoes em uma mesma carga.

\subsection{05\_dashboard\_misto/00\_dashboard\_defesa\_civil}

\textbf{Objetivo local.} simular painel territorial combinando municipios, hospitais e TIs\\

\textbf{Datasource(s).} dim\_municipios\_coordenadas, dim\_hospitais\_cnes, dim\_terras\_indigenas\\

\textbf{Padrão de carga.} dashboard\_composite\\

\textbf{Visualização Superset simulada.} tabela\_mapa\\

\textbf{Motivo científico.} Valida se o perfil suporta uma sessao analitica mais proxima do Superset. Objetivo local: simular painel territorial combinando municipios, hospitais e TIs\\

\textbf{Impacto esperado.} Exercita planejamento SQL, subconsultas, union, filtros e agregacoes em uma mesma carga.

\begin{lstlisting}[language=SQL]

WITH hospitais AS (
  SELECT
    codigo_ibge,
    COUNT(*) AS hospitais_total,
    SUM(CASE WHEN funciona_24h = 'SIM' THEN 1 ELSE 0 END) AS hospitais_24h,
    SUM(CASE WHEN dentro_terra_indigena = 1 THEN 1 ELSE 0 END) AS hospitais_ate_50km_ti
  FROM dim_hospitais_cnes
  GROUP BY codigo_ibge
),
terras AS (
  SELECT
    codigo_ibge_principal,
    COUNT(*) AS terras_indigenas_total,
    SUM(area_hectares) AS area_ti_hectares
  FROM dim_terras_indigenas
  GROUP BY codigo_ibge_principal
)
SELECT
  m.regiao_brasil,
  m.estado_sigla,
  COUNT(*) AS municipios_total,
  SUM(COALESCE(h.hospitais_total, 0)) AS hospitais_total,
  SUM(COALESCE(h.hospitais_24h, 0)) AS hospitais_24h,
  SUM(COALESCE(h.hospitais_ate_50km_ti, 0)) AS hospitais_ate_50km_ti,
  SUM(COALESCE(t.terras_indigenas_total, 0)) AS terras_indigenas_total,
  SUM(COALESCE(t.area_ti_hectares, 0)) AS area_ti_hectares
FROM dim_municipios_coordenadas m
LEFT JOIN hospitais h
  ON h.codigo_ibge = m.codigo_ibge
LEFT JOIN terras t
  ON t.codigo_ibge_principal = m.codigo_ibge
GROUP BY
  m.regiao_brasil,
  m.estado_sigla
ORDER BY hospitais_total DESC;

\end{lstlisting}

\subsection{05\_dashboard\_misto/01\_dashboard\_energia\_logistica}

\textbf{Objetivo local.} simular comparativo operacional entre energia e logistica\\

\textbf{Datasource(s).} dim\_ons\_linhas\_transmissao, dim\_ons\_subestacoes, dim\_rodovias\_federais\\

\textbf{Padrão de carga.} union\_aggregate\\

\textbf{Visualização Superset simulada.} tabela\_barras\\

\textbf{Motivo científico.} Valida se o perfil suporta uma sessao analitica mais proxima do Superset. Objetivo local: simular comparativo operacional entre energia e logistica\\

\textbf{Impacto esperado.} Exercita planejamento SQL, subconsultas, union, filtros e agregacoes em uma mesma carga.

\begin{lstlisting}[language=SQL]

SELECT
  'linhas_transmissao' AS camada,
  estado_sigla_estimado AS territorio,
  classe_criticidade_linha AS classe_risco,
  COUNT(*) AS ativos_total,
  SUM(extensao_km) AS medida_total
FROM dim_ons_linhas_transmissao
GROUP BY
  estado_sigla_estimado,
  classe_criticidade_linha
UNION ALL
SELECT
  'subestacoes' AS camada,
  'NAO INFORMADO' AS territorio,
  criticidade_subestacao AS classe_risco,
  COUNT(*) AS ativos_total,
  SUM(capacidade_mw) AS medida_total
FROM dim_ons_subestacoes
GROUP BY criticidade_subestacao
UNION ALL
SELECT
  'rodovias' AS camada,
  'NAO INFORMADO' AS territorio,
  criticidade_rodovia AS classe_risco,
  COUNT(*) AS ativos_total,
  COUNT(*) AS medida_total
FROM dim_rodovias_federais
GROUP BY criticidade_rodovia;

\end{lstlisting}

\subsection{05\_dashboard\_misto/02\_dashboard\_qualidade\_operacional}

\textbf{Objetivo local.} simular painel de governanca com qualidade geografica e relacionamentos\\

\textbf{Datasource(s).} dim\_hospitais\_cnes, dim\_terras\_indigenas, dim\_municipios\_coordenadas\\

\textbf{Padrão de carga.} quality\_dashboard\\

\textbf{Visualização Superset simulada.} indicadores\\

\textbf{Motivo científico.} Valida se o perfil suporta uma sessao analitica mais proxima do Superset. Objetivo local: simular painel de governanca com qualidade geografica e relacionamentos\\

\textbf{Impacto esperado.} Exercita planejamento SQL, subconsultas, union, filtros e agregacoes em uma mesma carga.

\begin{lstlisting}[language=SQL]

SELECT
  'hospitais_sem_municipio' AS regra,
  COUNT(*) AS avaliados,
  SUM(CASE WHEN codigo_ibge IS NULL OR codigo_ibge = 'NAO INFORMADO' THEN 1 ELSE 0 END) AS pendencias
FROM dim_hospitais_cnes
UNION ALL
SELECT
  'hospitais_fora_bbox' AS regra,
  COUNT(*) AS avaliados,
  SUM(CASE WHEN latitude < -35 OR latitude > 6 OR longitude < -75 OR longitude > -30 THEN 1 ELSE 0 END) AS pendencias
FROM dim_hospitais_cnes
UNION ALL
SELECT
  'terras_indigenas_sem_municipio_principal' AS regra,
  COUNT(*) AS avaliados,
  SUM(CASE WHEN codigo_ibge_principal IS NULL OR codigo_ibge_principal = 'NAO INFORMADO' THEN 1 ELSE 0 END) AS pendencias
FROM dim_terras_indigenas
UNION ALL
SELECT
  'municipios_sem_path_json' AS regra,
  COUNT(*) AS avaliados,
  SUM(CASE WHEN path_json IS NULL OR path_json = '' THEN 1 ELSE 0 END) AS pendencias
FROM dim_municipios_coordenadas;

\end{lstlisting}

\clearpage

\section{Fatos e Streaming}

\textbf{Conceito base.} Consultas pesadas diretamente sobre os maiores volumes de dados com agregacoes temporais.

\textbf{Hipótese testada.} A tabela fato e o gargalo e pode saturar o motor Scan e GroupBy do Druid sem filtros eficientes.

\textbf{Por que existe.} Simula carga baseada no historico sem camada semantica.

\textbf{Impacto esperado no Druid.} Pressao maxima no Historical Node.

\subsection{06\_fatos\_streaming/00\_count\_eventos\_total}

\textbf{Objetivo local.} medir latencia baseline em tabela fato (158K linhas)\\

\textbf{Datasource(s).} fato\_eventos\_incendio\\

\textbf{Padrão de carga.} count\_scan\\

\textbf{Visualização Superset simulada.} indicador\\

\textbf{Motivo científico.} Simula carga baseada no historico sem camada semantica. Objetivo local: medir latencia baseline em tabela fato (158K linhas)\\

\textbf{Impacto esperado.} Pressao maxima no Historical Node.

\begin{lstlisting}[language=SQL]

SELECT COUNT(*) AS total_eventos FROM fato_eventos_incendio;

\end{lstlisting}

\subsection{06\_fatos\_streaming/01\_eventos\_groupby\_bioma\_mes}

\textbf{Objetivo local.} agregacao bioma x mes - teste do motor GroupBy em fato com 158K linhas\\

\textbf{Datasource(s).} fato\_eventos\_incendio\\

\textbf{Padrão de carga.} groupby\_temporal\\

\textbf{Visualização Superset simulada.} heatmap\\

\textbf{Motivo científico.} Simula carga baseada no historico sem camada semantica. Objetivo local: agregacao bioma x mes - teste do motor GroupBy em fato com 158K linhas\\

\textbf{Impacto esperado.} Pressao maxima no Historical Node.

\begin{lstlisting}[language=SQL]

SELECT
  fk_bioma,
  TIME_FLOOR(__time, 'P1M') AS mes,
  COUNT(*) AS eventos,
  SUM(total_focos) AS focos_total,
  MAX(max_frp) AS pico_frp
FROM fato_eventos_incendio
GROUP BY 1, 2
ORDER BY pico_frp DESC;

\end{lstlisting}

\subsection{06\_fatos\_streaming/02\_eventos\_top\_frp\_ordenado}

\textbf{Objetivo local.} simular o padrao exato dos Virtual Datasets (GROUP BY + ORDER BY + LIMIT)\\

\textbf{Datasource(s).} fato\_eventos\_incendio\\

\textbf{Padrão de carga.} groupby\_ordered\_limit\\

\textbf{Visualização Superset simulada.} mapa\_scatter\\

\textbf{Motivo científico.} Simula carga baseada no historico sem camada semantica. Objetivo local: simular o padrao exato dos Virtual Datasets (GROUP BY + ORDER BY + LIMIT)\\

\textbf{Impacto esperado.} Pressao maxima no Historical Node.

\begin{lstlisting}[language=SQL]

SELECT
  fire_event_id, __time, latitude_centroide, longitude_centroide,
  fk_bioma, total_focos, soma_frp, max_frp, area_evento_km2
FROM fato_eventos_incendio
WHERE __time >= TIMESTAMP '2025-08-01' AND __time < TIMESTAMP '2025-11-01'
GROUP BY 1, 2, 3, 4, 5, 6, 7, 8, 9
ORDER BY max_frp DESC
LIMIT 50000;

\end{lstlisting}

\subsection{06\_fatos\_streaming/03\_focos\_scan\_temporal}

\textbf{Objetivo local.} medir motor Scan em tabela de 3.3M linhas com filtro temporal\\

\textbf{Datasource(s).} fato\_focos\_incendio\\

\textbf{Padrão de carga.} scan\_temporal\\

\textbf{Visualização Superset simulada.} tabela\\

\textbf{Motivo científico.} Simula carga baseada no historico sem camada semantica. Objetivo local: medir motor Scan em tabela de 3.3M linhas com filtro temporal\\

\textbf{Impacto esperado.} Pressao maxima no Historical Node.

\begin{lstlisting}[language=SQL]

SELECT
  __time, latitude, longitude, potencia_radiativa_fogo, risco_fogo_inpe,
  fk_bioma, fk_municipio
FROM fato_focos_incendio
WHERE __time >= TIMESTAMP '2025-09-01' AND __time < TIMESTAMP '2025-09-08'
LIMIT 10000;

\end{lstlisting}

\subsection{06\_fatos\_streaming/04\_focos\_lookup\_municipio}

\textbf{Objetivo local.} medir custo de LOOKUP em tabela de 3.3M linhas\\

\textbf{Datasource(s).} fato\_focos\_incendio\\

\textbf{Padrão de carga.} lookup\_heavy\\

\textbf{Visualização Superset simulada.} tabela\\

\textbf{Motivo científico.} Simula carga baseada no historico sem camada semantica. Objetivo local: medir custo de LOOKUP em tabela de 3.3M linhas\\

\textbf{Impacto esperado.} Pressao maxima no Historical Node.

\begin{lstlisting}[language=SQL]

SELECT
  fk_bioma,
  LOOKUP(fk_municipio, 'municipios_uf') AS uf,
  COUNT(*) AS focos,
  AVG(potencia_radiativa_fogo) AS media_frp,
  MAX(potencia_radiativa_fogo) AS max_frp
FROM fato_focos_incendio
WHERE __time >= TIMESTAMP '2025-08-01' AND __time < TIMESTAMP '2025-11-01'
GROUP BY 1, 2
ORDER BY focos DESC;

\end{lstlisting}

\clearpage

\section{Virtual Datasets}

\textbf{Conceito base.} Modelos multidimensionais com varios lookups simulando a camada semantica real.

\textbf{Hipótese testada.} Lookups sobre uma Fato grande nao deve derrubar a performance devido aos mapas compactados.

\textbf{Por que existe.} Representa o workload mais realista e pesado gerado pelos dashboards em producao.

\textbf{Impacto esperado no Druid.} Pressao distribuida entre Historical e Broker.

\subsection{07\_virtual\_datasets/00\_ds\_eventos\_geral}

\textbf{Objetivo local.} simular query real do dashboard de eventos geral (Virtual Dataset)\\

\textbf{Datasource(s).} fato\_eventos\_incendio\\

\textbf{Padrão de carga.} virtual\_dataset\\

\textbf{Visualização Superset simulada.} dashboard\_completo\\

\textbf{Motivo científico.} Representa o workload mais realista e pesado gerado pelos dashboards em producao. Objetivo local: simular query real do dashboard de eventos geral (Virtual Dataset)\\

\textbf{Impacto esperado.} Pressao distribuida entre Historical e Broker.

\begin{lstlisting}[language=SQL]

SELECT
  fire_event_id,
  __time,
  -- Geolocation & Basic info
  latitude_centroide,
  longitude_centroide,
  fk_bioma,
  
  -- Metrics
  total_focos,
  soma_frp,
  max_frp,
  area_evento_km2,
  duracao_dias,
  media_risco_fogo,
  
  -- Município Lookups
  fk_municipio_principal,
  LOOKUP(fk_municipio_principal, 'municipios_nome') AS municipio_nome,
  LOOKUP(fk_municipio_principal, 'municipios_uf') AS municipio_uf,
  LOOKUP(fk_municipio_principal, 'municipios_regiao') AS municipio_regiao,
  
  -- Saúde Metrics & Lookups
  indice_pressao_maxima,
  indice_pressao_saude,
  distancia_hospital_km,
  fk_hospital,
  LOOKUP(fk_hospital, 'hospitais_nome') AS hospital_nome,
  LOOKUP(fk_hospital, 'hospitais_tipo') AS hospital_tipo,
  LOOKUP(fk_hospital, 'hospitais_critico') AS hospital_critico,
  LOOKUP(fk_hospital, 'hospitais_24h') AS hospital_24h
FROM fato_eventos_incendio
WHERE 1=1
  AND __time >= TIMESTAMP '2025-08-01' AND __time < TIMESTAMP '2025-11-01'
GROUP BY 1, 2, 3, 4, 5, 6, 7, 8, 9, 10, 11, 12, 13, 14, 15, 16, 17, 18, 19, 20, 21, 22, 23
ORDER BY max_frp DESC
LIMIT 50000;

\end{lstlisting}

\subsection{07\_virtual\_datasets/01\_ds\_focos\_individuais}

\textbf{Objetivo local.} simular query real do dashboard de focos individuais (Virtual Dataset)\\

\textbf{Datasource(s).} fato\_focos\_incendio\\

\textbf{Padrão de carga.} virtual\_dataset\\

\textbf{Visualização Superset simulada.} dashboard\_completo\\

\textbf{Motivo científico.} Representa o workload mais realista e pesado gerado pelos dashboards em producao. Objetivo local: simular query real do dashboard de focos individuais (Virtual Dataset)\\

\textbf{Impacto esperado.} Pressao distribuida entre Historical e Broker.

\begin{lstlisting}[language=SQL]

SELECT
  __time,
  latitude,
  longitude,
  
  -- Metrics
  potencia_radiativa_fogo AS frp,
  risco_fogo_inpe,
  dias_sem_chuva,
  
  -- Dimensions & Lookups
  fk_bioma,
  fk_satelite_sensor,
  fk_municipio,
  LOOKUP(fk_municipio, 'municipios_nome') AS municipio_nome,
  LOOKUP(fk_municipio, 'municipios_uf') AS municipio_uf

FROM fato_focos_incendio
WHERE
  latitude IS NOT NULL
  AND longitude IS NOT NULL
  AND coordenada_valida = 1
  AND __time >= TIMESTAMP '2025-09-01' AND __time < TIMESTAMP '2025-09-08'
GROUP BY 1, 2, 3, 4, 5, 6, 7, 8, 9, 10, 11
ORDER BY frp DESC
LIMIT 50000;

\end{lstlisting}

\subsection{07\_virtual\_datasets/02\_ds\_eventos\_energia}

\textbf{Objetivo local.} simular query real do dashboard de eventos energia (Virtual Dataset)\\

\textbf{Datasource(s).} fato\_eventos\_incendio\\

\textbf{Padrão de carga.} virtual\_dataset\\

\textbf{Visualização Superset simulada.} dashboard\_completo\\

\textbf{Motivo científico.} Representa o workload mais realista e pesado gerado pelos dashboards em producao. Objetivo local: simular query real do dashboard de eventos energia (Virtual Dataset)\\

\textbf{Impacto esperado.} Pressao distribuida entre Historical e Broker.

\begin{lstlisting}[language=SQL]

SELECT
  fire_event_id,
  __time,
  total_focos,
  
  -- Linhas de Transmissão
  indice_ameaca_linha,
  distancia_linha_km,
  fk_linha_transmissao,
  LOOKUP(fk_linha_transmissao, 'linha_nome') AS linha_nome,
  CAST(LOOKUP(fk_linha_transmissao, 'linha_tensao') AS FLOAT) AS linha_tensao_kv,
  LOOKUP(fk_linha_transmissao, 'linha_uf') AS linha_uf,
  
  -- Subestações
  indice_ameaca_subestacao,
  distancia_subestacao_km,
  fk_subestacao,
  LOOKUP(fk_subestacao, 'subestacao_nome') AS subestacao_nome,
  CAST(LOOKUP(fk_subestacao, 'subestacao_capacidade') AS FLOAT) AS subestacao_capacidade_mw,
  CAST(LOOKUP(fk_subestacao, 'subestacao_betweenness') AS FLOAT) AS subestacao_betweenness,
  LOOKUP(fk_subestacao, 'subestacao_risco_cascata') AS subestacao_risco_cascata,
  
  -- Sistemas Isolados
  indice_ameaca_sistema_isolado,
  distancia_sistema_isolado_km,
  fk_sistema_isolado,
  LOOKUP(fk_sistema_isolado, 'sistema_isolado_nome') AS sistema_isolado_nome,
  LOOKUP(fk_sistema_isolado, 'sistema_isolado_uf') AS sistema_isolado_uf,
  max_frp

FROM fato_eventos_incendio
WHERE
  (indice_ameaca_linha > 0 
  OR indice_ameaca_subestacao > 0 
  OR indice_ameaca_sistema_isolado > 0)
  AND __time >= TIMESTAMP '2025-08-01' AND __time < TIMESTAMP '2025-11-01'
GROUP BY 1, 2, 3, 4, 5, 6, 7, 8, 9, 10, 11, 12, 13, 14, 15, 16, 17, 18, 19, 20, 21, 22
ORDER BY max_frp DESC
LIMIT 50000;

\end{lstlisting}

\subsection{07\_virtual\_datasets/03\_ds\_eventos\_indigenas}

\textbf{Objetivo local.} simular query real do dashboard de eventos indigenas (Virtual Dataset)\\

\textbf{Datasource(s).} fato\_eventos\_incendio\\

\textbf{Padrão de carga.} virtual\_dataset\\

\textbf{Visualização Superset simulada.} dashboard\_completo\\

\textbf{Motivo científico.} Representa o workload mais realista e pesado gerado pelos dashboards em producao. Objetivo local: simular query real do dashboard de eventos indigenas (Virtual Dataset)\\

\textbf{Impacto esperado.} Pressao distribuida entre Historical e Broker.

\begin{lstlisting}[language=SQL]

SELECT
  fire_event_id,
  __time,
  latitude_centroide,
  longitude_centroide,
  total_focos,
  soma_frp,
  area_evento_km2,
  duracao_dias,
  
  -- TI Metrics
  indice_pressao_terra_indigena,
  distancia_terra_indigena_km,
  
  -- TI Classification
  CASE
    WHEN fk_terra_indigena IS NOT NULL AND fk_terra_indigena != 'NAO INFORMADO'
    THEN 'DENTRO'
    ELSE 'ENTORNO'
  END AS tipo_ameaca,
  
  -- TI Lookups (Dentro)
  fk_terra_indigena,
  LOOKUP(fk_terra_indigena, 'ti_nome') AS ti_nome_dentro,
  LOOKUP(fk_terra_indigena, 'ti_etnia') AS ti_etnia_dentro,
  LOOKUP(fk_terra_indigena, 'ti_uf') AS ti_uf_dentro,
  
  -- TI Lookups (Proxima/Entorno)
  fk_terra_indigena_proxima,
  LOOKUP(fk_terra_indigena_proxima, 'ti_nome') AS ti_nome_proxima,
  LOOKUP(fk_terra_indigena_proxima, 'ti_etnia') AS ti_etnia_proxima,
  LOOKUP(fk_terra_indigena_proxima, 'ti_uf') AS ti_uf_proxima,

  -- Unified TI Lookups for easy filtering
  COALESCE(
    LOOKUP(fk_terra_indigena, 'ti_nome'),
    LOOKUP(fk_terra_indigena_proxima, 'ti_nome')
  ) AS ti_nome_unificado,
  COALESCE(
    LOOKUP(fk_terra_indigena, 'ti_etnia'),
    LOOKUP(fk_terra_indigena_proxima, 'ti_etnia')
  ) AS ti_etnia_unificada,
  max_frp
  
FROM fato_eventos_incendio
WHERE 
  ((fk_terra_indigena IS NOT NULL AND fk_terra_indigena != 'NAO INFORMADO')
  OR indice_pressao_terra_indigena > 0)
  AND __time >= TIMESTAMP '2025-08-01' AND __time < TIMESTAMP '2025-11-01'
GROUP BY 1, 2, 3, 4, 5, 6, 7, 8, 9, 10, 11, 12, 13, 14, 15, 16, 17, 18, 19, 20, 21, 22
ORDER BY max_frp DESC
LIMIT 50000;

\end{lstlisting}

\subsection{07\_virtual\_datasets/04\_ds\_eventos\_logistica}

\textbf{Objetivo local.} simular query real do dashboard de eventos logistica (Virtual Dataset)\\

\textbf{Datasource(s).} fato\_eventos\_incendio\\

\textbf{Padrão de carga.} virtual\_dataset\\

\textbf{Visualização Superset simulada.} dashboard\_completo\\

\textbf{Motivo científico.} Representa o workload mais realista e pesado gerado pelos dashboards em producao. Objetivo local: simular query real do dashboard de eventos logistica (Virtual Dataset)\\

\textbf{Impacto esperado.} Pressao distribuida entre Historical e Broker.

\begin{lstlisting}[language=SQL]

SELECT
  fire_event_id,
  __time,
  total_focos,
  
  -- Rodovias Metrics
  indice_ameaca_rodovia,
  distancia_rodovia_km,
  
  -- Rodovias Lookups
  fk_trecho_rodoviario,
  LOOKUP(fk_trecho_rodoviario, 'rodovia_codigo') AS rodovia_codigo,
  LOOKUP(fk_trecho_rodoviario, 'rodovia_criticidade') AS rodovia_criticidade,
  
  -- Município Relacionado
  fk_municipio_principal,
  LOOKUP(fk_municipio_principal, 'municipios_nome') AS municipio_nome,
  LOOKUP(fk_municipio_principal, 'municipios_uf') AS municipio_uf,
  max_frp

FROM fato_eventos_incendio
WHERE
  indice_ameaca_rodovia > 0
  AND fk_trecho_rodoviario IS NOT NULL
  AND fk_trecho_rodoviario != 'NAO_INFORMADO'
  AND __time >= TIMESTAMP '2025-08-01' AND __time < TIMESTAMP '2025-11-01'
GROUP BY 1, 2, 3, 4, 5, 6, 7, 8, 9, 10, 11, 12
ORDER BY max_frp DESC
LIMIT 50000;

\end{lstlisting}


\end{document}
